% conclusion.tex

\section{Conclusion}
\label{sec:conclusion}

We set out to understand why SF-GRU-family pedestrian crossing-intention
models fail to generalize across PIE and JAAD, and found more than one
answer was needed. Part of the difficulty was a silent, orthogonal data
bug in JAAD's default configuration, unrelated to cross-dataset transfer
itself but affecting any result built on the unaudited default; we
identify, fix, and re-confirm our findings hold under the correction.
The remaining, genuine cross-dataset collapse is not explained by camera
geometry, and the closely related phenomenon of split-sensitivity is
not explained by video-level train/test leakage -- both plausible,
natural hypotheses that we falsify directly rather than merely fail to
support. A domain-adversarial and contrastive training recipe, applied
to the fused representation and motivated by a modality-zeroing
ablation, produces a real, seed-confirmed improvement in both transfer
directions on SF-GRU, but leaves cross-dataset performance well below
in-domain levels; we consider this residual gap, not the improvement,
to be the more important number to report honestly. A
representation-separability probe, deliberately independent of the
domain classifier used during training, shows this recipe's benefit is
dissociated from the literal domain-invariance it is designed to
induce -- the representations that help most are, if anything,
\emph{more} separable by an external probe, not less -- suggesting its
practical effect is closer to a generic regularizer on the fused
representation than to representation alignment per se. Testing the
same recipe on two further, structurally distinct fusion architectures
supports this reading further and tempers any claim of a general-
purpose fix: on cross-attention fusion the improvement replicates only
partially, and on uncertainty-weighted fusion -- the strongest
unadapted baseline architecture we find -- the same recipe
significantly \emph{degrades} performance. We take from this that
domain-adversarial training's value here is real but architecture-
specific, not evidence that it achieves anything close to its nominal
objective of representational invariance regardless of architecture.

We draw two broader conclusions. First, cross-dataset evaluation should
be treated as standard practice for pedestrian crossing-intention
architectures, not an optional follow-up: every architecture we survey
in Section~\ref{sec:related} published since Gesnouin~et~al.'s original
finding~\cite{gesnouin2022} reports only in-domain results, and our own
in-domain numbers (Table~\ref{tab:baseline}) look entirely unremarkable
next to a near-chance cross-dataset collapse the same checkpoint
exhibits. In-domain performance on PIE or JAAD alone is not informative
about whether a proposed architecture has learned anything that
transfers. Second, we think auditing benchmark and pipeline details
that are easy to get silently wrong -- data-loader defaults, split
construction, the provenance of a supposedly domain-invariant
representation -- is undersupplied relative to the rate at which new
architectures are proposed for these benchmarks, and that the negative
results in this paper (ruled-out geometry, ruled-out leakage, a
dissociated invariance mechanism) are worth reporting even though none
of them, individually, is the paper's headline number.

\textbf{Limitations.} Our experiments center on SF-GRU's stacked fusion;
we additionally test our headline recipe on two further, structurally
distinct architectures (box-anchored cross-attention,
Section~\ref{sec:results-crossattn}, and uncertainty-weighted fusion,
Section~\ref{sec:results-uncertainty}). The probe dissociation --
higher separability wherever the recipe changes task behavior at all --
replicates on both, but the \emph{fix itself} does not: cross-attention
shows a partial, one-direction improvement, and uncertainty-weighted
fusion shows the recipe significantly harming its (strongest, of any
architecture we test) unadapted baseline. We therefore have reasonable
evidence the dissociation mechanism is not an SF-GRU-specific artifact,
but active evidence \emph{against} the fix being architecture-general,
not merely an absence of evidence for it. We did not test transformer-
or Mamba-based architectures at all, so whether our findings extend to
substantially larger or differently-inductive-biased fusion mechanisms
remains open, and a practitioner adopting this recipe on a new
architecture should re-validate both directions rather than assume
transfer. We evaluate on two datasets; both are
pedestrian-crossing-specific,
relatively small by modern standards, and share enough protocol
lineage (both from the same research group) that a third, more
independently sourced dataset (e.g.\ TITAN) would strengthen any claim
that our findings reflect a general property of this task rather than
a property specific to the PIE/JAAD pair. Our camera-geometry
falsification relies on a researched, not measured, camera calibration
for JAAD, since JAAD does not publish one; we mitigate this with a
sensitivity sweep over plausible parameter ranges but cannot rule out
that a parameter setting outside our swept range behaves differently.
Finally, our synthesized JAAD speed modality is a coarse optical-flow
proxy, not a calibrated measurement, and its negative result should be
read as evidence about this particular proxy rather than as a general
claim that a speed signal cannot help JAAD cross-dataset transfer.

We release our corrected-labeling configuration, falsification
experiments, and probe methodology in the hope that they are directly
reusable by future work evaluating new architectures on PIE and JAAD,
independent of whether that work adopts our specific training recipe.
